\definecolor{lightlightgrey}{rgb}{0.9,0.9,0.9}
\definecolor{col1}{RGB}{255, 230, 230}
\definecolor{col2}{rgb}{0.9372, 0.5804, 0.6196}
\definecolor{col3}{RGB}{233, 255, 245}
%%%%%%%
\begin{document}
\begin{table*}[t]
\centering
\caption{
\textbf{Ablation study.}
We report top-1 accuracy}
\vspace{-2.5mm}
    \begin{subtable}[t]{0.43\textwidth}
        \renewcommand{\arraystretch}{1.1}
        \centering
        \caption{Learning ``bursty'' prompts by our BPT methods outperforms prior arts VPT and SPT.
        Refer to Fig.~\ref{fig:splashy-figure} for ``bursty'' distributions.}
        \resizebox{!}{0.165\textwidth}{
        \begin{tabular}{lccrcc}
        \toprule
        methods & init. & tuning? & \#params & IN-1K & CUB-200\\
        \midrule
        \rowcolor{lightlightgrey} VPT~\cite{jia2022visual}           & -      & -            & 7.68  & 63.71 & 42.15\\
        \rowcolor{lightlightgrey} SPT~\cite{wang2024revisiting}      & -      & -            & 7.68  & 69.98 & 71.15\\
        BPT-fWhiten                         & whiten & \ding{55}    & 7.68  & 72.09 & 77.48\\
        BPT-tWhiten                         & whiten & $\checkmark$ & 66.66 & \textbf{72.37} & \textbf{78.54}\\
        \cellcolor{col2}{BPT-bilinear}  & random & $\checkmark$ & \textbf{6.51}  & 72.15 & 77.86\\
        \bottomrule
        \end{tabular}
        }
        \label{tab:prompt-strategy}
    \end{subtable}
    \hfill
    \begin{subtable}[t]{0.25\textwidth}
        \renewcommand{\arraystretch}{1.1}
        \centering
        \caption{\textbf{Prompt length}. Increasing length further improves accuracy.}
        \resizebox{!}{0.285\textwidth}{
        \begin{tabular}{cccc}
        \toprule
        length & params & IN-1K & CUB-200\\
        \midrule
        64                       & 6.24 & 71.92 & 76.91\\
        81                       & 6.37 & 71.94 & 77.60\\
        \cellcolor{col2}{100} & 6.51 & 72.15 & 77.86\\
        144                      & 6.84 & 72.39 & 78.36\\
        196                      & 7.23 & \textbf{72.48} & \textbf{78.87}\\
        \bottomrule
        \end{tabular}
        }
        \label{tab:prompt-length}
    \end{subtable}
    \hfill
    \begin{subtable}[t]{0.28\textwidth}
        \renewcommand{\arraystretch}{1.1}
        \centering
        \caption{
        \textbf{Prompt width}. BPT-bilinear learns more compact prompts with higher accuracy.
        }
        \resizebox{!}{0.255\textwidth}{
        \begin{tabular}{cccc}
        \toprule
        width & params & IN-1K & CUB-200\\
        \midrule
        \rowcolor{lightlightgrey} SPT~\cite{wang2024revisiting}      & 7.68  & 69.98 & 71.15\\
        25                      & 2.17 & 72.06 & 76.54\\
        50                      & 4.34 & 72.08 & 76.95\\
        \cellcolor{col2}{75} & 6.51 & 72.15 & 77.86\\
        100                     & 8.68 & \textbf{72.18} & \textbf{78.59}\\
        \bottomrule
        \end{tabular}
        }
        \label{tab:prompt-width}
    \end{subtable}
\vspace{-1mm}
\label{tab:ablation}
\end{table*}
\setlength{\tabcolsep}{0.6em}
\begin{table}[t]
\centering
\caption{\small
\textbf{Analysis of BPT-Deep design choices.}
Extending from shallow to deep, BPT achieves significant improvements.
Recall our BPT (-bilinear) allows flexible design of prompt to control the total number of learned parameters.
We study how to design the bilinear prompt and where to insert prompt in a Transformer network.
Applying prompt tuning in all Transformer blocks achieves the best performance, but inserting prompts at the top four or six blocks is still competitive with much fewer learned parameter.
}
\vspace{-3mm}
\scalebox{0.8}{
    \begin{tabular}{l cccc}
\toprule
        variant & \#blocks & params ($\times$$10^{-2}$M) & IN-1K & CUB-200\\
\midrule
        \rowcolor{lightlightgrey}BPT-Shallow  & 1  & 6.51  & 72.15 & 77.86\\
        BPT-Deep  & 12 & 49.08 & \textbf{74.73} & \textbf{82.78}\\
        BPT-Deep  & 6 & 26.04 & 74.62 & 82.33\\
        BPT-Deep  & \cellcolor{col2}4 & 18.36 & 74.25 & 82.00\\
        BPT-Deep  & 3 & 14.52 & 74.07 & 80.86\\
    \midrule
        \rowcolor{lightlightgrey} SPT-Deep~\cite{wang2024revisiting} & 12 & 18.43 & 73.34 & 80.13\\
\bottomrule
\end{tabular}
}
\vspace{-5mm}
\label{tab:deep-variant}
\end{table}
\setlength{\tabcolsep}{0.34em}
\begin{table}[t]
\centering
\caption{\small
\textbf{Prompt tuning (pt) \textit{vs}. backbone fune-tuning (ft)} with a ViT-B  supervised pretrained on ImageNet-21K.
For each method, we report its parameters (in $10^{-2}$M) and the mean accuracy over three random runs on five FGVC datasets.
Our BPT outperforms the compared pt methods, ft methods, and even the full fine-tuning approach, with far fewer learned parameters.
}
\vspace{-3mm}
\scalebox{0.8}{
    \begin{tabular}{l cclrc}
    \toprule
        Methods & pt & ft & \textcolor{gray}{venue\&year} & \#param. &  mean Acc (\%)\\
        \midrule
        \rowcolor{lightlightgrey} Full fine-tuning       & ~ & $\checkmark$ & \textcolor{gray}{---} & 8.60$\times 10^9$  & 88.54\\
        \rowcolor{lightlightgrey} Linear probing                      & ~ & ~            & \textcolor{gray}{---} & 0.0   & 79.32\\
        \rowcolor{lightlightgrey} Bias~\cite{zhai2019large}           & ~ & $\checkmark$ & \textcolor{gray}{arXiv 2019} & 10.71 & 88.41\\
        \rowcolor{lightlightgrey} Adapter~\cite{houlsby2019parameter} & ~ & $\checkmark$ & \textcolor{gray}{ICML 2019} & 23.70 & 85.66\\
        \rowcolor{lightlightgrey} MP~\cite{gao2023tuning}             & ~ & $\checkmark$ & \textcolor{gray}{ICCV 2023} & 112.68 & 89.38\\
        \rowcolor{lightlightgrey} LoRA~\cite{hu2021lora}              & ~ & $\checkmark$ & \textcolor{gray}{ICLR 2022} & 39.46 & 89.85\\
        \rowcolor{lightlightgrey} SSF~\cite{lian2022scaling}          & ~ & $\checkmark$ & \textcolor{gray}{NeurIPS 2022}  & 21.69 & 90.72\\
        \rowcolor{lightlightgrey} SNF~\cite{wang2023adapting}         & ~ & $\checkmark$ & \textcolor{gray}{CVPR 2023}  & 7.68  & 90.74\\
        \rowcolor{col1} VPT-Shallow~\cite{jia2022visual}        & $\checkmark$ & & \textcolor{gray}{ECCV 2022}  & 7.68 & 84.62\\
        \rowcolor{col1} SPT-Shallow~\cite{wang2024revisiting}   & $\checkmark$ & & \textcolor{gray}{ICML 2024}  & 7.68 & 90.10\\
        \rowcolor{col1} {\bf BPT-Shallow}    & $\checkmark$ & & \textcolor{gray}{\bf ours} & \underline{6.51} & \underline{90.54}\\
        \rowcolor{col3} VPT-Deep~\cite{jia2022visual}            & $\checkmark$ & & \textcolor{gray}{ECCV 2022}     & 66.98 & 89.11\\
        \rowcolor{col3} E$^2$VPT~\cite{cheng2023e2vpt}           & $\checkmark$ & & \textcolor{gray}{ICCV 2023}     & 37.98 & 89.22\\
        \rowcolor{col3} SPT-Deep~\cite{wang2024revisiting}       & $\checkmark$ & & \textcolor{gray}{ICML 2024}     & 18.43 & 91.40\\
        \rowcolor{col3} {\bf BPT-Deep}                    & $\checkmark$ & & \textcolor{gray}{\bf ours}    & \textbf{18.36} & \textbf{91.72}\\
    \bottomrule
    \end{tabular}
}
\vspace{-3mm}
\label{tab:sota_sup}
\end{table}
\end{document}
